\documentclass{article}
% Chemistry: structural formulas (chemfig), equations (mhchem, chemformula),
% units (siunitx), hazard statements and reaction schemes (chemmacros),
% molecular orbital diagrams (modiagram), and a periodic-table style cell.
\usepackage[version=4]{mhchem}
\usepackage{chemfig}
\usepackage{chemformula}
\usepackage{chemmacros}
\chemsetup{modules={reactions,redox,phases,isotopes,nomenclature}}
\usepackage{siunitx}
\usepackage{modiagram}
\usepackage[margin=2cm]{geometry}

\begin{document}
\section{Equations}
\ce{2H2 + O2 -> 2H2O}, \quad \ce{CO2 + C <=> 2CO}, \quad \ce{^{227}_{90}Th+},
\quad \ce{SO4^2- + Ba^2+ -> BaSO4 v}, \quad \ce{Hg^2+ ->[I-] HgI2 ->[I-] [Hg^{II}I4]^2-}.

With chemformula: \ch{Fe2O3 + 3 CO -> 2 Fe + 3 CO2}.
A concentration of \SI{0.1}{\mole\per\litre} at \SI{298.15}{\kelvin}; $\Delta H = \SI{-285.8}{\kilo\joule\per\mole}$.

\begin{reaction}
  CH4 + 2 O2 -> CO2 + 2 H2O
\end{reaction}

Oxidation states: \ox{+2,Fe}, \ox{-2,O}; \iupac{2-chloro\-butane}; \pH{} and \pKa{}.

\section{Structures}
\chemfig{*6((-OH)=-=(-COOH)-=-)}
\qquad
\chemfig{H_3C-[:30]C(=[:90]O)-[:-30]O-[:30]CH_3}
\qquad
\chemfig{[:-30]*6(=-=(-[:30]N(-[:90]H)-[:-30]H)-=-)}

\schemestart
  \chemfig{C(-[2]H)(-[4]H)(-[6]H)-OH}
  \arrow{->[\ce{H+}][\SI{100}{\celsius}]}
  \chemfig{C(-[2]H)(-[4]H)=O}
  \+ \ce{H2O}
\schemestop

\section{Molecular orbitals}
\begin{MOdiagram}[labels]
  \atom{left}{1s={;pair}}
  \atom{right}{1s={;up}}
  \molecule{1sMO={.75;pair}}
\end{MOdiagram}

\section{Names}
\iupac{1,3-di\-chloro\-benzene} and \ch{C6H4Cl2}.
\end{document}
